% !TeX root = ../../Book1.tex
% 纲要标题：### 7.4 交错群的单性
% 纲要位置：计划纲要.md，第 320 行。

\section{交错群的单性}
\label{b1:sec:0704}

% 纲要内容（仅作写作计划，尚非教材正文）：
%
% * 三轮换生成与正规子群分析
% * n≥5 时交错群的单性
% * 与一般方程不可根式求解的后续联系
%
% ---
%

\subsection{三轮换生成与正规子群}
\label{b1:subsec:0704:01}

一个非平凡正规子群一旦包含某种置换，就必须包含它的所有群内共轭。因此，证明交错群单性的关键是找到足够简单且共轭丰富的元素。三轮换最适合这个任务：它是移动字母数最少的非单位偶置换。

\begin{lemma}\label{b1:lem:three-cycles-generate}
三轮换生成 \(A_n\)。当 \(n\ge5\) 时，所有三轮换在 \(A_n\) 中共轭。
\end{lemma}
\begin{proof}
三轮换生成 \(A_n\) 已在\secref{b1:sec:0302}证明；这里只需补充群内共轭性。

任意两个三轮换 \(\alpha,\beta\) 在 \(S_n\) 中共轭，设 \(u\alpha u^{-1}=\beta\)。若 \(u\) 为偶置换，已经得到所需共轭。若 \(u\) 为奇置换，因 \(n\ge5\)，可在 \(\beta\) 的支集之外选两个字母构成换位 \(t\)。它与 \(\beta\) 交换，且 \(tu\) 为偶置换，于是 \((tu)\alpha(tu)^{-1}=\beta\)。
\end{proof}
因此只要证明非平凡正规子群包含一个三轮换，就能推出它等于整个 \(A_n\)。下面的交换子方法把任意非单位元改造成移动较少字母的元素。

\begin{definition}\label{b1:def:permutation-support}
置换 \(\sigma\) 的\term[zhichi]{支集}[support]是
\(\operatorname{supp}(\sigma)=\{i\mid\sigma(i)\ne i\}\)。其大小称为此处的支持度。
\end{definition}
若 \(N\trianglelefteq A_n\)、\(\sigma\in N\)、\(\tau\in A_n\)，则
\[
[\sigma,\tau]=\sigma\tau\sigma^{-1}\tau^{-1}\in N.
\]
如果 \(\tau\) 是三轮换，两个乘积因子的支集各有三个字母，故交换子的支持度至多为六。还必须确保交换子非平凡；下述证明在每次使用时都核对这一点。

\subsection{交错群的单性}
\label{b1:subsec:0704:02}

\begin{theorem}\label{b1:thm:alternating-simple}
当 \(n\ge5\) 时，交错群 \(A_n\) 是非阿贝尔单群。
\end{theorem}
\begin{proof}
设 \(1\ne N\trianglelefteq A_n\)，在 \(N\setminus\{1\}\) 中取支持度最小的元素 \(\sigma\)，记该支持度为 \(m\)。偶置换不可能恰好移动一个或两个字母，所以 \(m\ge3\)。

先证 \(m\le6\)。若 \(m>6\)，选 \(a\) 满足 \(\sigma(a)=b\ne a\)，再选不同于 \(a,b\) 的两个字母 \(c,d\)，令 \(\tau=(acd)\)。共轭 \(\sigma\tau\sigma^{-1}\) 的支集包含 \(b\)，而 \(\tau\) 的支集不包含 \(b\)，故二者不同，\([\sigma,\tau]\ne1\)。这个交换子属于 \(N\)，支持度至多六，违背 \(m\) 的最小性。

若 \(m=3\)，则 \(\sigma\) 已是三轮换。其余可能的偶置换轮换类型只有下面几种；字母均表示互不相同的点。
\begin{enumerate}
\item 若 \(m=4\)，必有 \(\sigma=(ab)(cd)\)。选支集外的字母 \(e\)（此处用 \(n\ge5\)），令 \(\tau=(cde)\)。于是
\(\sigma\tau\sigma^{-1}=(dce)=\tau^{-1}\)，所以
\([\sigma,\tau]=\tau^{-2}=\tau\) 是 \(N\) 中的三轮换。这也表明支持度最小元不能实际具有 \(m=4\)。
\item 若 \(m=5\)，偶性迫使 \(\sigma=(abcde)\)。取 \(\tau=(abc)\)，则共轭为 \((bcd)\)。二者支集不同，所以交换子非平凡；其支集包含于 \(\{a,b,c,d\}\)，违背最小性。
\item 若 \(m=6\) 且 \(\sigma=(abc)(def)\)，取 \(\tau=(abd)\)。共轭为 \((bce)\)，所以交换子非平凡且支集包含于 \(\{a,b,c,d,e\}\)，仍与最小性矛盾。
\item 若 \(m=6\) 且 \(\sigma=(abcd)(ef)\)，取 \(\tau=(abc)\)。共轭为 \((bcd)\)，产生一个非平凡且支持度至多四的交换子，再次矛盾。
\end{enumerate}
这些情形穷尽了支持度四至六的偶置换：六轮换、三个换位，以及三轮换乘一个换位都是奇置换；含固定点的轮换型应按实际支持度计入前面的情形。因此 \(N\) 含有三轮换。由前一引理及正规性，\(N\) 包含所有三轮换，故 \(N=A_n\)。最后，\((123)\) 与 \((124)\) 不交换，说明 \(A_n\) 非阿贝尔。
\end{proof}
\term[jiaocuo-qun-danxing]{交错群的单性}[simplicity of the alternating group]在 \(n=5\) 就成立；证明没有使用多于五个字母去处理 \(A_5\) 的双换位。低阶例外也应同时记住：\(A_3\cong C_3\) 是阿贝尔单群，\(A_4\) 有真正规子群 \(V_4\)，因而不单。

\subsection{交错群与根式不可解性}
\label{b1:subsec:0704:03}

\begin{corollary}\label{b1:cor:symmetric-nonsolvable}
当 \(n\ge5\) 时，\(A_n'=A_n\)，因而 \(A_n\) 及 \(S_n\) 均不可解。对 \(n\le4\)，\(S_n\) 可解。
\end{corollary}
\begin{proof}
导出子群正规。由于 \(A_n\) 非阿贝尔，其导出子群不平凡；单性迫使它等于 \(A_n\)。所以导出列不可能下降到单位群，\(A_n\) 不可解，而可解性传给子群表明 \(S_n\) 也不可解。低阶情形在\secref{b1:sec:0702}已经算出；\(S_1,S_2\) 则直接阿贝尔。
\end{proof}
这为一般方程的根式不可解性准备了群论部分。到本卷 Galois 理论中，我们将建立“用有限次开根表达根”与适当 Galois 群的可解性之间的联系，并证明相应一般多项式的 Galois 群是对称群。两项桥梁建立之后，\(S_n\) 在 \(n\ge5\) 时不可解才转化为一般五次及更高次方程没有统一根式公式。

\ProseParagraphBreak
当前结果并不说“每个五次方程都不可根式求解”。例如 \(x^5-2=0\) 的一个根本来就是 \(\sqrt[5]{2}\)，其余根再乘五次单位根即可表达。决定具体方程的是它自己的 Galois 群，而不是次数一个数字。这里先明确逻辑分工：本章完成可解群及非阿贝尔单群的事实，域扩张与根式之间的对应留到本卷后续章节证明。
